\documentclass[10pt,a4paper]{amsart}
\usepackage[margin=19mm]{geometry}
\usepackage{amsmath,amssymb,mathtools}
\usepackage[scr=rsfs,cal=euler]{mathalfa}
\usepackage{xfrac}
\usepackage[unicode,hidelinks,bookmarks=false]{hyperref}
\newcommand{\ie}{i.e.\ }
\newcommand{\cf}{cf.\ }
\newcommand{\resp}{respectively\ }
\newcommand{\Subsection}{\subsection}
\providecommand{\nobreakdash}{\nobreak\hbox{-}}
\setlength{\emergencystretch}{2em}
\multlinegap0pt
\usepackage{bm}
\usepackage{bbm}
\usepackage{dsfont}
\usepackage{xspace}
\usepackage{cancel}
\usepackage{mathtools}
\usepackage{amsthm}
\makeatletter
\@ifundefined{lemma}{\newtheorem{lemma}{Lemma}}{}
\@ifundefined{proposition}{\newtheorem{proposition}{Proposition}}{}
\makeatother
\newcommand{\FF}{{\mathcal F}}
\newcommand{\FPSIPF}[1]{FPS-IP(#1)}
\newcommand{\FPSIP}{FPS-IP\xspace}
\newcommand{\PP}{{\mathcal P}}
\title[Capacitated power dominating set problem: a solution approac]{Capacitated power dominating set problem: a solution approach based on forbidden propagation sets}
\author{Mauro Lucci, Diego Delle Donne, Mariana Escalante}
\date{Source: arXiv:2605.18533v1 (CC BY 4.0)}
\begin{document}
\maketitle

\noindent\emph{Source and licence.} Capacitated power dominating set problem: a solution approach based on forbidden propagation sets. Version arXiv:2605.18533v1 (CC BY 4.0), distributed under the Creative Commons Attribution 4.0 International licence, as recorded in the arXiv licence metadata for this exact version and shown on the abstract page https://arxiv.org/abs/2605.18533v1. Full rights notice: licenses/arxiv-2605.18533-CC-BY-4.0.txt.\par\smallskip

\noindent\textit{Editorial scope.} Counted unit: the Proposition whose source label is \texttt{prop.FPS.IP} in the article (source0.tex lines 1257--1259, source label \texttt{prop.FPS.IP}), delivered together with the article's own complete original proof of it (source0.tex lines 1261--1295, one \texttt{proof} environment of 35 lines). No mathematics is written, repaired or re-derived: statement and proof are the article's own text line for line. Required context retained uncounted: lemma.DAG (lines 495--498); FPSIP.R1 (lines 1213--1219); FPSIP.R2 (lines 1213--1219); FPSIP.R3 (lines 1213--1219); lemma.fps.constraint.generation (lines 1232--1238). Every definition, display and label of those lines is retained unchanged. Omitted from this package and not redistributed: 1--494, 499--1212, 1220--1231, 1239--1256, 1260--1260, 1296--1790. Nothing inside a retained excerpt is omitted or altered: every retained body is delivered exactly as the article's lines are written, so no omission receipt of any kind accompanies this package. Citations: the retained excerpts carry no citation command at all, so no bibliography entry of the article is transcribed and no reference is redistributed. Reference adaptation: none was needed and none was made; every reference inside the delivered unit resolves inside it, so no unresolved reference or citation is produced. Printed slips kept literal: no printed word, symbol, hypothesis or number of the counted statement or of its proof was altered. Layout and class: the article's own class is \texttt{article} with packages including authblk, amsmath, mathrsfs, mathtools, amssymb, xspace, amsthm, enumitem, algorithm, algpseudocode, tikz, subcaption, microtype, hyperref; the wrapper is the lane's pinned \texttt{amsart} editorial wrapper at 10pt on A4, which restates the theorem-like environments the retained text needs and adds the article's own macro definitions that the retained text needs (\texttt{\textbackslash FF}, \texttt{\textbackslash FPSIP}, \texttt{\textbackslash FPSIPF}, \texttt{\textbackslash PP}), with extra wrapper packages (bm, bbm, dsfont, xspace, cancel, mathtools). The delivered numbering is the unit's own. Assets: no retained span contains a figure environment, a graphics-inclusion command, a PDF-page inclusion, a table or a TikZ picture; no image file is copied into this package, the package declares no asset dependency and no image file is redistributed with it. Licence: version arXiv:2605.18533v1 of this article is distributed under the Creative Commons Attribution 4.0 International licence, as recorded in the arXiv licence metadata for this exact version and shown on the abstract page https://arxiv.org/abs/2605.18533v1. A full rights notice is delivered at licenses/arxiv-2605.18533-CC-BY-4.0.txt. Characters: every retained excerpt is pure ASCII, so the delivered engine, XeLaTeX with its default Latin Modern text font, reports no missing character.
\begin{lemma}\label{lemma.DAG}
Let $\rho$ be a $k$-capacitated function.
If $R \subseteq A_P$ is the set of propagations applied by rule (PR) in a proper calculation of $M(\rho)$, then the precedence digraph $D_R$ contains no cycles.
\end{lemma}

\begin{align}
\text{(\FPSIP)}\ \min\  & \sum_{v \in V} s_v \label{FPSIP.ObjFun}\\
s.t.\  & s_v + \sum_{\mathclap{\substack{u \in N(v):\\\delta(u) \leq k}}} s_u + \sum_{\mathclap{\substack{u \in N(v):\\\delta(u) > k}}} w_{uv} + \sum_{\mathclap{u \in N(v) \cap V_P}} y_{uv} \geq 1 & v \in V \label{FPSIP.R1}\\
& \sum_{u \in N(v)} w_{vu} \leq ks_v & v \in V: \delta(v) > k \label{FPSIP.R2}\\
& \sum_{(u,v) \in F} y_{uv} \leq |F|-1 & F \in \FF(G,V_P) \label{FPSIP.R3}\\
& s \in \{0,1\}^V,\ w \in \{0,1\}^{A_D},\ y \in \{0,1\}^{A_P}. \notag
\end{align}

\begin{lemma}\label{lemma.fps.constraint.generation}
Let $\tilde \FF \subseteq \FF(G,V_P)$ and $(\tilde s, \tilde w, \tilde y)$ be a feasible solution of \FPSIPF{$\tilde \FF$}.
Define $\rho: S_{\rho} \to \PP(V)$ as $S_{\rho} = \{v \in V: \tilde s_v = 1\}$ and, for each $u \in S_{\rho}$, $\rho(u) = N(u)$ if $\delta(u) \leq k$ and $\rho(u) = \{v \in N(u): \tilde w_{uv} = 1\}$ if $\delta(u) > k$.

If $\rho$ is not a $k$-capacitated power dominating function, then the set $F = \{(u,v) \in A_P: \tilde y_{uv} = 1\}$ is an FPS.
Moreover, the precedence digraph $D_F$ contains a  cycle whose vertices all belong to $V \setminus M(\rho)$. 
\end{lemma}

\begin{proposition}\label{prop.FPS.IP}
\FPSIP is a correct formulation for CPDS.
\end{proposition}

\begin{proof}
Let $\rho$ be a $k$-capacitated power dominating function. 
Thus, there exists a finite sequence of applications of rules (DR) and (PR) such that $M(\rho) = V$, which we may assume to be proper.
We first show that there exists a feasible solution of \FPSIP that encodes $\rho$.

Consider the vector $(\tilde s, \tilde w, \tilde y) \in \{0,1\}^{V \times A_D \times A_P}$ defined as: $\tilde s_v = 1$ if and only if $v \in S_{\rho}$, for all $v \in V$; $\tilde w_{uv} = 1$ if and only if the rule (DR) is applied at $u$ to monitor $v$, for all $(u,v) \in A_D$; $\tilde y_{uv} = 1$ if and only if the rule (PR) is applied at $u$ to monitor $v$, for all $(u,v) \in A_P$.
We will see that $(\tilde s, \tilde w, \tilde y)$ verifies the constraints of \FPSIP and therefore is a feasible solution.

To prove the validity of constraints \eqref{FPSIP.R1}, let $v \in V$. 
As $M(\rho) = V$, $v$ must be monitored by rule (DR) or (PR).
If $v$ monitors itself through rule (DR), then $\tilde s_v = 1$.
But if a neighbor $u$ monitors $v$ through rule (DR), then $\tilde s_u = 1$ when $\delta(u) \leq k$, and $\tilde w_{uv} = 1$ otherwise.
Finally, if some neighbor $u \in V_P$ monitors $v$ through rule (PR), then $\tilde y_{uv} = 1$.
Thus, the left-hand side is at least 1.

For constraints \eqref{FPSIP.R2}, let $v \in V$ with $\delta(v) > k$. 
Since $\rho$ is a $k$-capacitated power dominating function, we have $|\rho(v)| \leq k$. Hence, $v$ monitors at most $k$ neighbors through rule (DR), which implies $\sum_{u \in N(v)} \tilde w_{vu} \leq k$. 
Additionally, if $v \notin S_{\rho}$, rule (DR) cannot be applied at all, implying $\tilde w_{vu} = 0$ for all $u \in N(v)$. 

For constraints \eqref{FPSIP.R3}, let $R \subseteq A_P$ be the set of propagations applied by rule (PR) in the calculation of $M(\rho)$. By Lemma \ref{lemma.DAG}, $D_R$ contains no  cycles, which guarantees that there is no FPS $F$ such that $\tilde y_{uv} = 1$ for all $(u,v) \in F$.

\medskip

Now, given a feasible solution $(\tilde s, \tilde w, \tilde y)$ of \FPSIP.
Define $S_{\rho} = \{v \in V: \tilde s_v = 1\}$ and $\rho: S_{\rho} \to \PP(V)$ by setting, for each $u \in S_{\rho}$, $\rho(u) = N(u)$ if $\delta(u) \leq k$ and $\rho(u) = \{v \in N(u): \tilde w_{uv} = 1\}$ if $\delta(u) > k$.

Suppose that $\rho$ is not a $k$--capacitated power dominating function.
By Lemma \ref{lemma.fps.constraint.generation}, the precedence digraph $D_F$, where $F = \{(u,v) \in A_P: \tilde y_{uv} = 1\}$, contains a  cycle.
Hence, $F$ is an FPS.
Therefore, $F$ contains a minimal FPS $F'$, and the FPS constraint associated with $F'$ would be violated by $(\tilde s, \tilde w, \tilde y)$, contradicting the feasibility of the solution.

\medskip

Finally, in the proposed constructions for both directions, the cardinality of $S_{\rho}$ is equal to the objective value of the solution $(\tilde s, \tilde w, \tilde y)$.
\end{proof}

\end{document}
